% !PRINTED 483
\weberpage{483}
Ist $-i\omega > \sqrt2$, so ist auch $f_1^2(\omega) > \sqrt2$, denn $f_1(\omega)^2$ geht,
während $-i\omega$ von $\sqrt2$ bis $\infty$ geht, ebenfalls, und zwar stets
wachsend\footnote{Aus § 54, (11) folgt nämlich durch die Substitution $(\omega, 1-\frac1{\omega})$:
\[ d\,f_1(\omega)^8 = -\frac{\pi i}{3}\vartheta_{10}^4\left[f(\omega)^8+f_1(\omega)^8\right]d\omega. \]},
von $\sqrt2$ bis $\infty$, und folglich ist $A_1$ negativ für
$r = 0, 2, 4$.
Der erste Faktor $A_1+1$ verschwindet, wie aus dem schon
bekannten Resultat (24) hervorgeht, für $r = 4$; daher verschwindet
der dritte Faktor $A_1^3+2A_1^2+A_1+4$ für $r = 0$, während
$A_1+2$ für $r = 2$ verschwindet.
Wir erhalten also nach (25):
\weq{27}{f_1(\sqrt{-10})^2 = \sqrt2\,y, \quad y - \frac1y = 1,}
\weq{28}{f_1(\sqrt{-22})^2 = \sqrt2\,y, \quad y - \frac1y = 2,}
\weq{29}{f_1(\sqrt{-26})^2 = \sqrt2\,y, \quad y^6-2y^5-2y^4+2y^2-2y-1=0.}

\begin{center}\rule{3cm}{0.4pt}\end{center}

IV. Als Beispiel für eine andere Art der Verwendung der
\gesperrt{Schlaeflischen} Modulargleichungen, die zu der sonst schwer zugänglichen
 Klasseninvariante $f(\sqrt{-41})$ führt, möge das Folgende
dienen.
Wir setzen:
\weq{30}{2n\omega^2 + 2r\omega + n = 0,}
\weq{31}{\omega = \frac{-r+\sqrt{-m}}{2n}, \quad m = 2n^2-r^2.}
Es kann hierin $r$ jede Zahl bedeuten, die $m$ positiv macht,
die aber mit $n$ keinen Teiler gemeinschaftlich hat [weil sonst
(30) imprimitiv ist]. Es wird dann nach § 34:
\weq{32}{f_2\left(\frac{\omega}{n}\right) = f_1[2(n\omega+r)] = \frac{e^{-\frac{r\pi i}{12}}\sqrt2}{f_2(n\omega)},}
\[ f_1[2(n\omega+r)] = f_1(r+\sqrt{-m}) = e^{-\frac{r\pi i}{24}}x, \]
%CONT>>
